\subsection{A finite-ring equation determines the local constraint}

The last step of \cite[Theorem~5]{Sanz2010Quantum} passes from one finite
ring to all sufficiently long rings. The interaction is fixed throughout.

\begin{lemma}[Local annihilation of every virtual boundary]
    \label{lem:local_weighted_boundary_annihilation}
    \lean{MPSTensor.weighted_boundary_eq_zero_of_periodicLocalInteractionES_eq_zero}
    \leanok
    \uses{lem:cyclic_mpv_active_boundary}
    Let the complementary tuples of words in $A_1,\ldots,A_b$ span the
    product matrix algebra at length $N-L$, with $L\le N$, and let $c_j\in\C$.
    If a local translate $h_i$ annihilates $\sum_jc_jV^{(N)}(A_j)$, then
    \begin{align}
        c_jh\Gamma_L^{A_j}(X)=0
        \qquad\text{for every }j\text{ and every virtual matrix }X.
        \label{eq:local_weighted_boundary_annihilation}
    \end{align}
    The operator $h$ may be any linear operator, and the coefficients may vanish.
\end{lemma}
\begin{proof}\leanok
    \uses{lem:cyclic_mpv_active_boundary}
    Compose the joint boundary map with $h$, extend its vanishing from the
    spanning tuples to the whole product algebra, and evaluate on a boundary
    supported in one block.
\end{proof}

\begin{lemma}[The entire active MPS space is in the interaction kernel]
    \label{lem:finite_volume_local_kernel_inclusion}
    \lean{MPSTensor.boundary_eq_zero_of_periodicInteractionHamiltonianES_eq_zero,
        MPSTensor.groundSpaceES_le_ker_of_periodicInteractionHamiltonianES_eq_zero}
    \leanok
    \uses{lem:local_weighted_boundary_annihilation, def:periodic-fixed-interaction}
    Suppose $N>0$, $L\le N$, $h\ge0$, and every $c_j\ne0$. If the length-$N-L$
    word tuples span the product matrix algebra and
    $H_N\sum_jc_jV^{(N)}(A_j)=0$, then
    \begin{align}
        h\Gamma_L^{A_j}(X) & =0\quad\text{for every }X,
                           & G_L(A_j)                   & \subseteq\ker h\quad\text{for every }j.
        \label{eq:finite_volume_local_kernel_inclusion}
    \end{align}
\end{lemma}
\begin{proof}\leanok
    \uses{lem:local_weighted_boundary_annihilation}
    Positivity separates the finite Hamiltonian equation into local equations.
    Use the translate at site zero and
    Lemma~\ref{lem:local_weighted_boundary_annihilation}, then cancel $c_j$.
\end{proof}

\begin{lemma}[A fixed local kernel annihilates every permitted periodic volume]
    \label{lem:local_mps_kernel_all_volumes}
    \lean{MPSTensor.periodicLocalInteractionES_mpvState_eq_zero_of_groundSpaceES_le_ker,
        MPSTensor.periodicInteractionHamiltonianES_mpvState_eq_zero_of_groundSpaceES_le_ker}
    \leanok
    \uses{lem:cyclic_mpv_active_boundary, def:periodic-fixed-interaction}
    Let $h$ be linear on $L$ sites and suppose $G_L(A)\subseteq\ker h$.
    For every $M\ge L$, every cyclic translate of $h$ annihilates the actual
    periodic vector $V^{(M)}(A)$. Consequently $H_MV^{(M)}(A)=0$.
    No positivity or complementary-span condition at $M$ is needed.
\end{lemma}
\begin{proof}\leanok
    \uses{lem:cyclic_mpv_active_boundary}
    Every complementary configuration supplies a virtual boundary matrix for
    the active block. The local-kernel inclusion kills each such fiber, and
    summing the translated equations gives the Hamiltonian equation.
\end{proof}

\begin{lemma}[Propagation from a finite ring to every permitted volume]
    \label{lem:finite_ring_constraint_all_volumes}
    \lean{MPSTensor.periodicInteractionHamiltonianES_mpvState_eq_zero_all_volumes}
    \leanok
    \discussion{10280}
    \uses{lem:finite_volume_local_kernel_inclusion, lem:local_mps_kernel_all_volumes}
    Let $N>0$, $L\le N$, and $h\ge0$. If the words of length $N-L$ in $A$
    span its full matrix algebra and $H_NV^{(N)}(A)=0$, then
    \begin{align}
        H_MV^{(M)}(A)=0\qquad\text{for every }M\ge L.
        \label{eq:finite_ring_constraint_all_volumes}
    \end{align}
    The same tensor $A$ and interaction $h$ work at every volume. Nonvanishing
    of these periodic vectors and uniqueness of the ground state, needed for
    the full dichotomy, are separate assertions.
\end{lemma}
\begin{proof}\leanok
    \uses{lem:finite_volume_local_kernel_inclusion, lem:local_mps_kernel_all_volumes}
    Apply Lemma~\ref{lem:finite_volume_local_kernel_inclusion} to the one-block
    family with coefficient one. Then apply
    Lemma~\ref{lem:local_mps_kernel_all_volumes}.
\end{proof}
